% !TEX root = main.tex

\renewcommand{\arraystretch}{1.5}
Consideramos el caso de 2 $e^-$ con  $l=1$. Las posibles 
configuraciones electrónicas serán aquellas que 
obedezcan el Principio de Exclusión de Pauli. Para 
encontrarlas, vamos a considerar las posibles 
configuraciones de espín  en función del 
número de duplas $(m_{l}^{(1)},m_{l}^{(2)})$ con $m_{l}^{(1)}=m_{l}^{(2)}$.
En este caso, al tener solo 2 $e^-$, el número máximo 
de duplas iguales es 1, pero en general dependerá del valor de  $l$. 
%NOTE:(Se puede buscar cuál es el nro. máximo de duplas que se pueden formar en función de $l$ y  $N$).
Una vez determinamos el número máximo de duplas 
permitido $\mathcal{D}_{M}$, es fácil ver que van a 
existir configuraciones permitidas con valor de
duplas $\mathcal{D}\in [0,\mathcal{D}_{M}]$. A cada valor de dupla $\mathcal{D}$ 
le es asociado un conjunto de posibles configuraciones de 
espín.\par 
Para el caso particular de $2$  $e^-$ con  $l=1$, 
tenemos  $\mathcal{D}_{M}=1$ y las \textbf{<<reglas>>} o 
posibles configuraciones de espín para 
cada $\mathcal{D}$ son:\par 
\underline{\textbf{Reglas}}
\begin{align*}
  \mathcal{D}=1&\longrightarrow \begin{cases}
    1\times M_{s}=0:\quad (+,-)
  \end{cases}\\
    \mathcal{D}=0&\longrightarrow \begin{cases}
      1\times M_{s}=1:\quad (+,+)\\
      2\times M_{s}=0:\quad (+,-)
    \end{cases}
\end{align*}
Es quizás importante remarcar, en el caso de $\mathcal{D}=1$, 
que las configuraciones $(+,-)$ y $(-,+)$ no son distintas 
y sólo se cuenta una de ellas. Esto es porque en este
caso  $m_{l}^{(1)}=m_{l}^{(2)}$, y como el resto de 
números cuánticos también son iguales, los 
$e^-$ son equivalentes, y el estado 
de esta configuración debe ser la combinación antisimétrica de las configuraciones mencionadas 
 \begin{align}
  \ket{\psi} = \frac{1}{\sqrt{2}}\left(\ket{n,n;l,l;m_{l},m_{l};+,-}-\ket{n,n;l,l;m_{l},m_{l};-,+} \right) 
\end{align}
es decir ambas configuraciones refieren al mismo estado. 
Estamos empleando la siguiente notación para 
especificar una configuración 
\begin{align*}
  \ket{\psi} = \ket{n_{1},n_{2};l_{1},l_{2};m_{l}^{(1)},m_{l}^{(2)};m_{s}^{(1)},m_{s}^{(2)}}
\end{align*}
Mientras que en el caso $\mathcal{D}=0$, al tener $m_{l}$ distintos, 
los dos electrones son inequivalentes, y por lo tanto 
cuando $M_{S}=0$ tenemos dos posibles estados diferentes 
\begin{align*}
  \ket{\psi_{1}} &= \frac{1}{\sqrt{2}}\left(\ket{n,n;l,l;m_{l},m_{l}';+,-}-\ket{n,n;l,l;m_{l}',m_{l};-,+}\right)\\
 \ket{\psi_{2}} &= \frac{1}{\sqrt{2}}\left(\ket{n,n;l,l;m_{l},m_{l}';-,+}-\ket{n,n;l,l;m_{l}',m_{l};+,-}\right)\\
  \ket{\psi_{1}}&\ne \ket{\psi_{2}}
\end{align*}
Establecidas las reglas, procedemos a obtener los términos según el valor de $M_{L}$:
\begin{align*}
 \begin{array}{c | c | c | c}
  M_{L} & \text{ Particiones } & \text{ Duplas } & \text{ Términos }\\
  \hline 
  M_{L} = 2 & 2 = 1 + 1 & \mathcal{D} = 1,\; k = 0 & 1 \times (M_{S} = 0) \implies \ce{^{1}D}(5)\\
  \hline 
  M_{L} = 1 & 1 = 1 + 0 & \mathcal{D} = 0,\; k = 2 & \begin{cases}
    1 \times (M_{S} = 1) \implies \ce{^{3}P}(9)\\
    2 \times (M_{S} = 0) \implies \cancel{\ce{^{1}D}(5)},\ \cancel{\ce{^{3}P}(9)}
    \end{cases}\\
    \hline 
  M_{L} = 0 & 0 = 1 - 1 & \mathcal{D} = 0,\; k = 2 & \begin{cases}
    1 \times (M_{S} = 1) \implies \cancel{\ce{^{3}P}(9)}\\
    2 \times (M_{S} = 0) \implies \cancel{\ce{^{1}D}(5)},\ \cancel{\ce{^{3}P}(9)}
    \end{cases}\\
  & 0 = 0 + 0 & \mathcal{D} = 1,\; k = 0 & 1 \times (M_{S} = 0) \implies \ce{^{1}S}(1)\\
  \hline
  \end{array}
  \end{align*}
Una de las ventajas de este método es que, a la hora 
de determinar las posibles combinaciones de $m_{l}$ que 
suman $M_{L}$ basta con partir del $M_{L}^{(\text{max})}$ máximo que 
se obtiene trivalmente de maximizar todos los $m_{l}$ 
(respetando el principio de Pauli), y luego para 
$M_{L}^{(\text{max})}-1$ las combinaciones posibles se obtienen 
de restar $1$ a los diferentes $m_{l}$ que aparecen 
en la configuración de $M_{L}^{(\text{max})}$.\par 
Por otro lado, al disminuir $M_{L}$ y contar los 
estados, tenemos que considerar que para una 
configuración L-S dada, también van a 
existir instancias con  $M_{L}$ y $M_{S}$ menores, 
por lo que en cada $M_{L}$ debemos de asegurarnos 
de no perder estos estados ni contarlos como 
configuraciones nuevas o distintas. Específicamente a tener en cuenta en la tabla anterior es que, en cada $M_{L}$ hemos contabilizado y posteriormente tachado los términos que ya han surgido en $M_{L}$ mayores, y que, por lo tanto, no son nuevos.\par

Para $M_{L}<0$ las configuraciones serán 
idénticas, simplemente cambiando el signo de las $m_{l}$. 
Entonces, los términos espectrales que obtenemos son:
\begin{align*}
  \Aboxed{\text{T.E.:}\quad &\ce{^{1}D}(5), \quad \ce{^{3}P}(9),\quad \ce{^{1}S}(1)}\\ 
  \text{ Degeneración total: }&g_{T} = 5 + 9 + 1 = 15 = \binom{6}{2}\quad \boxed{\checkmark}
\end{align*}
